\section{Introduction}
\label{sec:intro}

EEG foundation models tokenize a recording before they model it, and the
choice of what a token carries decides what the model can learn cheaply and
what it can learn only by accident. Across the current generation the
choice varies in how much of the signal survives tokenization --- magnitude
spectrograms \citep{yang2023biot,pradeepkumar2026tfm}, raw patches
\citep{wang2025cbramod,wang2024eegpt}, per-bin spectral amplitude and phase
\citep{jiang2024labram} --- but not in one respect: every token describes
one frequency band, or one patch of the broadband waveform, on its own. No
token type represents the \emph{relation} between two bands, and no
pretraining objective supervises one.

That relation is not an exotic quantity. Phase--amplitude coupling, the
locking of a fast band's envelope to the phase of a slow rhythm, has a
standard estimator \citep{tort2010measuring}, a physiology documented for
two decades \citep{canolty2006high,canolty2010functional}, and direct
clinical relevance to precisely the tasks foundation models are evaluated
on: it is stronger in the seizure onset zone than in unaffected tissue,
interictally as well as ictally \citep{amiri2016phase}, and periodic
epileptiform discharges are, by definition, fast activity riding a slow
rhythm. A representation from which such a quantity could in principle be
decoded is not yet one that encodes it. Magnitude tokenizers make it hard
to reach, because re-deriving phase from overlapping magnitude frames is a
phase-retrieval problem \citep{griffin1984signal} that nothing in the
pipeline performs; waveform tokenizers keep it reachable but give the
encoder no incentive to extract it, and probing studies find that the
resulting embeddings under-represent exactly the oscillatory structure
such relations live in \citep{kommineni2026aperiodic}.

This paper asks what a clinical EEG model gains from cross-frequency
structure, and finds that the answer depends on where the structure is
put. We build a small model whose tokenizer decomposes each electrode into
learned frequency bands and, optionally, measures the coupling between
them per patch and hands it to the encoder as a first-class token; we then
evaluate the two ingredients separately, on twelve clinical corpora, with
an evaluation designed so that either answer can be negative.

\paragraph{Contributions.}
\begin{itemize}[leftmargin=*,itemsep=2pt,topsep=2pt]
\item \textbf{A 1.6M-parameter model that is competitive with 25--69M
  foundation models on twelve clinical corpora, and far ahead on seizure
  detection.} Trained from scratch, it wins or ties the strongest
  reproduced baseline on seven corpora and is within a few points on the
  rest (Section~\ref{sec:main}). Architectural controls locate the margin:
  removing the learned band decomposition or the factorized tri-axial
  attention over it collapses seizure detection to the level of the large
  baselines. The result is about inductive bias in the clinical label
  regime, not about scale.
\item \textbf{Cross-frequency coupling as token content, with its scope
  measured.} Each band contributes a waveform token and an interaction
  token built from the band's measured coupling to slower bands; the
  interaction token is invariant to the phase reference of the analysis,
  preserves band energy exactly, and is gated so that the uncoupled model
  is an exactly reachable configuration (Section~\ref{sec:tokenizer}). With
  the encoder fixed, the token is decisive on epileptiform event
  classification ($+0.15$ Cohen's $\kappa$ on TUEV over three seeds,
  concentrated on periodic discharges) and neutral on the ten other corpora
  (Section~\ref{sec:coupling-results}). Transplanted into CBraMod's encoder
  in place of its own patch embedding, it improves event classification and
  ties on seizure detection, and in both cases switching the coupling off
  inside the transplant gives the gain back (Section~\ref{sec:transplant}).
\item \textbf{A pretraining study reported to convergence rather than
  assumed.} We construct the masked-reconstruction objective that
  supervises coupling directly, compare it against band-normalized
  amplitude reconstruction \citep{yu2026understanding} with lr-matched
  frozen probes, label-fraction curves, and clean-transfer controls, and
  report the pretrained model under one fixed finetuning recipe: it helps
  where labels are scarce and hurts where they are plentiful
  (Section~\ref{sec:pretrain-results}).
\end{itemize}

Two boundaries are stated up front. The coupling token is a mechanism for
a specific kind of label --- a cross-frequency morphology --- and we claim
no more than that for it; the general competitiveness of the model rests
on its band-decomposed tokens and factorized encoder. And sleep staging,
where coupling is abundant \citep{amiri2016phase} but the model is not
best, is reported as a boundary condition rather than folded into the
average.
